\documentclass[11pt]{article}
\usepackage{amsmath,amssymb,amsthm}
\usepackage{mathtools}
\usepackage{tikz}
\usepackage{enumitem}
\usepackage[margin=1in]{geometry}

\newtheorem{theorem}{Theorem}
\newtheorem*{theorem*}{Theorem}
\newtheorem{lemma}{Lemma}
\newtheorem{definition}{Definition}
\newtheorem*{definition*}{Definition}
\newtheorem{proposition}{Proposition}
\newtheorem{assumption}{Assumption}
\newtheorem{remark}{Remark}



\title{Math 106 -- Notes \\ Week 8: February 24, 2026}
\author{Ethan Levien}
\date{}

\begin{document}
\maketitle



\section*{Monte Carlo simulation and variance reduction}

The goal of Monte Carlo simulation is to approximate expectations with simulations. Consider the expectation 
\begin{equation}
\mu(f) = {\mathbb E}_{\mu}[f(X)]
\end{equation}
 where $f(X) \in L^2(\mu)$. We want to find an estimator $\hat{f}$ of $\mu(f)$ which has a small mean-squared error
 \begin{equation}
 {\rm MSE}_{\hat{f}} = {\mathbb E}[(\hat{f}-\mu(f))^2]
 \end{equation}
 This breaks down into Bias and variance: 
 \begin{equation}
  {\rm MSE}_{\hat{f}} = {\rm var}(\hat{f}) + \big({\mathbb E}[\hat{f}]-\mu(f)\big)^2
 \end{equation}
 A consistent estimator is one for which ${\rm MSE}_{\hat{f}}  \to 0$ as $m \to \infty$. 
 
 In the simplest case we have iid samples $X_1,\dots,X_m \sim \mu$ which are used to obtain the approximation 
\begin{equation}
\hat{f}_m \approx \frac{1}{m}\sum_{i=1}^mf(X_i)
\end{equation}
This is an unbiased estimator whose variance is ${\rm var}(\hat{f}_m) = {\rm var}_{\mu}(f(X))/m$. Thus, 
\begin{equation}
  {\rm MSE}_{\hat{f}_m}= {\rm var}(\hat{f}_m)< \epsilon^2 \implies  m > \frac{{\rm var}_{\mu}(f(X))}{\epsilon^2}
\end{equation}
This idea of variance reduction is to find a new variable $\tilde{X}$ which has the same expectation but a smaller variance. 

\subsection*{Importance sampling (Section 4.3)}
Importance sampling is one of the main variance reduction techniques. It essentially amounts to performing a change of variables in the expectation in order to sample from a lower variance distribution. Let us suppose that $\mu$ has a density and (abusing notation) write $\mu(x)$ for this density. Now let $q(x)$ be another density such that $\mu/q>0$ for all $x$. 
We can write
\begin{equation}
{\mathbb E}_{\mu}[f(X)] = {\mathbb E}_{q}\left[f(X)L(X)\right],\quad L(X) = \frac{\mu(X)}{q(X)}
\end{equation}
This suggests the estimator 
\begin{equation}
\hat{f}_{m,q} \approx \frac{1}{m}\sum_{i=1}^mf(\tilde{X}_i)L(\tilde{X}_i)
\end{equation}
where $\tilde{X}_i$ are iid samples from $q$. This works to our advantage if 
\begin{equation}
{\rm var}_q(f(X)L(X)) < {\rm var}_{\mu}(f(X)). 
\end{equation}


\section*{Motiving Example: Estimating the time to a rare event}

 Now let's see how we can perform a change of measure to perform importance sampling on the path space with a concrete example. Consider again an OU process, 
\begin{equation}\label{eq:sdeexample}
dX_t = -\theta X_t\,dt + \sqrt{2D}\,dW_t,
\end{equation}
and suppose we are interested in an event that occurs with (state-dependent) rate $\lambda(X_t)$. We assume $\lambda$ is an increasing function of $x$. Conditional on ${\mathcal F}_t^X$, the survival probability is
\begin{equation}
{\mathbb P}(\tau>t\mid {\mathcal F}_t^X)
=
\exp\!\left(-\int_0^t \lambda(X_s)\,ds\right),
\end{equation}
and therefore
\begin{equation}
{\mathbb P}(\tau>t)
=
{\mathbb E}_{\mu}\!\left[
\exp\!\left(-\int_0^t \lambda(X_s)\,ds\right)
\right],
\end{equation}
where $\mu$ denotes the law of the path $\{X_s\}_{0\le s\le t}$.

We would like to sample from another SDE where $\tau$ happens sooner (and hopefully we don't pay a price in the simulation cost). 

\subsection*{Detour: Radon--Nikodym derivatives}

To make use of importance sampling in the context of SDEs we need to understand when a change of measure on \emph{path space} is valid. In finite dimensions we often write densities $\mu(x)$ and $q(x)$, but for diffusion processes the relevant objects are probability measures on a space of paths. In this setting, the analogue of having a density is \emph{absolute continuity}.

Let $(\Omega,\mathcal F,(\mathcal F_t))$ be a filtered measurable space, and let $\mu$ and $q$ be probability measures on $(\Omega,\mathcal F)$. 
\begin{definition}
We say that $\mu$ is absolutely continuous with respect to $q$ if 
\begin{equation}
\mu \ll q
\end{equation}
if $q(A)=0$ implies $\mu(A)=0$ for all $A\in\mathcal F$. 
\end{definition}
In this case, the Radon--Nikodym theorem guarantees the existence of a nonnegative observable $L$ such that
\begin{equation}
\mu(A) = \int_A L(x)q(dx) \qquad \text{for all } A\in\mathcal F.
\end{equation}
We then write $L = d\mu/dq$ and call it the Radon--Nikodym derivative (or likelihood ratio). Equivalently, for a random variable $f$,
\begin{equation}
\mu(f) = {\mathbb E}_{\mu}[f] = {\mathbb E}_{q}[f\,L] = q(fL).
\end{equation}
In applications to SDEs we typically consider $\mathcal F_t$ (the information up to time $t$) and require absolute continuity \emph{locally in time},
\begin{equation}
\mu\big|_{\mathcal F_t} \ll q\big|_{\mathcal F_t},
\end{equation}
so that a likelihood ratio $L_t = d\mu/dq|_{\mathcal F_t}$ exists as an adapted process.

\subsection*{What goes wrong when changing measures on path space}
As an example, consider two OU processes on $[0,T]$:
\begin{align}
dX_t &= -\theta X_t\,dt +\sqrt{2D}\,dW_t^{(1)},\\
dY_t &= -\theta Y_t\,dt + \sqrt{2\tilde{D}}\,dW_t^{(2)},
\end{align}
where $W^{(1)}$ and $W^{(2)}$ are Wiener processes (possibly on different probability spaces). Let $\mu$ and $q$ denote the laws on the path space of $X_t$ and $Y_t$ respectively. 

A quick way to see what goes wrong is to consider the case $D\neq \tilde D$. 
As I showed in class, the quadratic variations of $X_t$ and $Y_t$ are 
\begin{equation}
[X,X]_T = 2DT
\qquad
[Y,Y]_T = 2\tilde D T
\end{equation}
respectively where the equalities should be interpreted in the almost sure sense (meaning they occur with probability 1 under $\mu$ and $q$ respectively). 
Define the measurable event in path space
\begin{equation}
A \coloneqq \{\omega\in C([0,T]) : [\omega]_T = 2DT\}.
\end{equation}
Here $\omega$ could be replaced by $X$ or $Y$ but we write it this way to indicate this is an event that can be considered under either $\mu$ or $q$. 
Then
\begin{equation}
\mu(A)=1,
\qquad
q(A)=0 \quad \text{if } \tilde D\neq D,
\end{equation}
Therefore $\mu$ is \emph{not} absolutely continuous with respect to $q$. In fact the measures are mutually singular,
meaning there exists an event of probability one under $\mu$ and probability zero under $q$ (and vice versa). Consequently, the Radon--Nikodym derivative $d\mu/dq$ does not exist on path space when $D\neq \tilde D$, and there is no likelihood ratio that can convert $\mu$--expectations into $q$--expectations for general path functionals.


Note that even when $D\neq \tilde D$, the finite-dimensional distributions at fixed times $\{t_1,\dots,t_n\}$ are mutually absolutely continuous (both are nondegenerate Gaussians on ${\mathbb R}^n$). The problem, however, is that the likelihood ratio blows up as $dt \to 0$. 
To see this, consider the finite-dimensional distribution of $\sqrt{2D}W_t$ at times $0,dt,2dt,\dots,ndt$: 
\begin{equation}
\rho_D(w_1,\dots,w_n) = \prod_{i=1}^n (4 \pi D dt)^{-1/2}e^{-(\Delta w_i)^2/(4D dt)} = (4 \pi D dt)^{-n/2}e^{-1/(4D dt) \sum_{i=1}^n(\Delta w_i)^2} 
\end{equation}
The ratio
\begin{equation}
\frac{\rho_D(w_1,\dots,w_n)}{\rho_{\tilde{D}}(w_1,\dots,w_n)} = \left(\tilde{D}/D \right)^{n/2} e^{(\tilde{D}^{-1}-D^{-1})\sum_{i=1}^n(\Delta w_i)^2/(4 dt)}.
\end{equation}
Although $dt$ has canceled in the ratio, there is still a factor of $1/dt$ in the exponent which causes it to blow-up as $dt \to 0$. 
 
 
 \subsection*{Changing the drift and Girsanov's Theorem}
Since we cannot change the diffusion coefficient, we need to adjust the drift term. Returning to the context of our modeling example -- that is, the problem of estimating the CDF of $\tau$ -- let's explore this possibility.  
To this end, let $\{u_t\}_{0\le t\le T}$ be an adapted control (often $u_t=u(X_t)$). Define a new probability measure $q$ on ${\mathcal F}_t$ under which $X$ solves
\begin{equation}
dX_t
=
\Big(-\theta X_t + \sqrt{2D}\,u_t\Big)\,dt
+
\sqrt{2D}\,dW_t^{(q)}.
\end{equation}
Let $q$ denote the induced law of the path $\{X_s\}_{0\le s\le t}$.

% Grid: t_i = i\,dt, i=0,1,\dots,n with dt = T/n. Write x_i \approx X_{t_i}.
% Let u_{i-1} denote u_{t_{i-1}} (or u(x_{i-1}) if feedback).
% NOTE: in the q-transition below, the OU drift term should be +\theta x_{i-1}dt, not +\theta dt.


Using the Euler--Maruyama approximation (we will eventually take the limit $dt \to 0$, so the order of the approximation doesn't matter), we get
\begin{equation}
p_q(x_i\mid x_{i-1})
=
\frac{1}{\sqrt{4\pi D\,dt}}
\exp\!\left(
-\frac{\big(x_i - x_{i-1} + \theta x_{i-1}dt - \sqrt{2D}\,u_{i-1}dt\big)^2}{4D\,dt}
\right),
\end{equation}
while under the original OU law $\mu$,
\begin{equation}
p_{\mu}(x_i\mid x_{i-1})
=
\frac{1}{\sqrt{4\pi D\,dt}}
\exp\!\left(
-\frac{\big(x_i - x_{i-1} + \theta x_{i-1}dt \big)^2}{4D\,dt}
\right).
\end{equation}


Since the prefactors match, their ratio is
\begin{align}
\frac{\rho_\mu(x_1,\dots,x_n\mid x_0)}{\rho_q(x_1,\dots,x_n\mid x_0)}
&=
\exp\!\left\{
-\frac{1}{4D\,dt}\sum_{i=1}^n
\Big[
\big(\Delta x_i+\theta x_{i-1}dt\big)^2
-
\big(\Delta x_i+\theta x_{i-1}dt-\sqrt{2D}\,u_{i-1}dt\big)^2
\Big]
\right\},
\end{align}
Using $(a)^2-(a-b)^2 = 2ab-b^2$ with
\(
a=\Delta x_i+\theta x_{i-1}dt
\)
and
\(
b=\sqrt{2D}\,u_{i-1}dt,
\)
we obtain
\begin{align}
\log\frac{\rho_\mu}{\rho_q}
&=
-\frac{1}{4D\,dt}\sum_{i=1}^n
\left[
2\big(\Delta x_i+\theta x_{i-1}dt\big)\big(\sqrt{2D}\,u_{i-1}dt\big)
-\big(\sqrt{2D}\,u_{i-1}dt\big)^2
\right] \\
&=
-\sum_{i=1}^n u_{i-1}\,\frac{\Delta x_i+\theta x_{i-1}dt}{\sqrt{2D}}
+\frac12\sum_{i=1}^n u_{i-1}^2\,dt.
\end{align}
Now, under $q$ the Euler increment can be written as
\begin{equation}
\Delta x_i+\theta x_{i-1}dt
=
\sqrt{2D}\,u_{i-1}dt + \sqrt{2D}\,\Delta W_i^{(q)}
\end{equation}
so
\begin{align}
\log\frac{\rho_\mu}{\rho_q}
&=
-\sum_{i=1}^n u_{i-1}\,\big(u_{i-1}dt+\Delta W_i^{(q)}\big)
+\frac12\sum_{i=1}^n u_{i-1}^2\,dt \\
&=
-\sum_{i=1}^n u_{i-1}\,\Delta W_i^{(q)}
-\frac12\sum_{i=1}^n u_{i-1}^2\,dt.
\end{align}
Therefore the finite-dimensional likelihood ratio is
\begin{equation}
\frac{\rho_\mu(x_1,\dots,x_n\mid x_0)}{\rho_q(x_1,\dots,x_n\mid x_0)}
=
\exp\!\left(
-\sum_{i=1}^n u_{i-1}\,\Delta W_i^{(q)}
-\frac12\sum_{i=1}^n u_{i-1}^2\,dt
\right).
\end{equation}
So the continuous-time Radon--Nikodym derivative is
\begin{equation}
L_t \coloneqq \frac{d\mu}{dq}\Big|_{\mathcal F_t}
=
\exp\!\left(
-\int_0^t u_s\,dW_s^{(q)}
-\frac12\int_0^t u_s^2\,ds
\right).
\end{equation}

Combining everything: 
\begin{equation}\label{eq:muimportance}
\mu(f) =  q(fL_t) = {\mathbb E}_{q}\!\left[
\exp\!\left(-\int_0^t \lambda(X_s)\,ds-\int_0^t u_s\,dW_s^{(q)}
-\frac12\int_0^t u_s^2\,ds
\right)
\right] 
\end{equation}
This is an application of \emph{Girsanov's theorem}, which tells us how to change measure of Wiener paths. 

\subsection*{Precise statement of Girsanov's Theorem 1}

\begin{theorem}[Special case of T9.2]
Consider the It\^o process 
\begin{equation}
d\tilde{W}_t = \phi(t,\omega)dt + dW_t, \quad \tilde{W}_0 = 0
\end{equation}
defined on some probability space $(\Omega, {\mathcal F},{\mathbb P})$. 
Assume 
\begin{equation}
{\mathbb E}\left[\exp\left(\frac{1}{2}\int_0^T|\phi(s,\omega)|^2ds \right) \right]< \infty
\end{equation}
Now let $\tilde{\mathbb P}$ be another probability measure on $(\Omega, {\mathcal F})$ defined by 
\begin{equation}
\frac{d\tilde{\mathbb P}}{d{\mathbb P}}\Big|_{{\mathcal F}_t} =  \exp\left(- \int_0^t\phi(s,\omega)dW_s - \frac{1}{2}\int_0^t\phi(s,\omega)^2ds \right)
\end{equation}
Then $\tilde{W}_t$ is a Wiener process under $\tilde{\mathbb P}$. 
\end{theorem}

\begin{remark}
To connect to our notation, the distribution of $W_t$ under ${\mathbb P}$ is $\mu$ -- $\mu$ is the law of a Wiener process. ${\mathbb P}$ induces a distribution over $\tilde{W}_t$ which is our $q$. $q$ is not the law of a Wiener process because of the drfit term, but the change of measure converts it back to one. 
\end{remark}


\section*{Feynman-Kac formula (Section 8.6)}
In the previous example we encountered the path expectation 
\begin{equation}
{\mathbb E}_{\mu}\!\left[
\exp\!\left(-\int_0^t \lambda(X_s)\,ds\right)
\right]
\end{equation}
The challenge of sampling from these expectations is pervasive. Because this is not the expectation of $X_t$ at a fixed time conditional on the initial point, it does not obey the usual backward equation. It does however have a semigroup structure.
\begin{equation}
v(x,t)  = {\mathbb E}_{\mu}\!\left[f(X_t)
\exp\!\left(-\int_0^t \lambda(X_s)\,ds\right)|X_0=x
\right]. 
\end{equation}
Note that if we define $Z_t = Z_0+\int_0^t\lambda(X_s)\,ds$, we have a new system of SDEs
\begin{align}
dZ_t &= \lambda(X_t)dt\\
dX_t &= b(X_t)dt + \sigma(X_t)dW_t
\end{align}
with this definition of $Z_t$. 
This has the generator 
\begin{equation}
({\mathcal A}g)(x,z) = ({\mathcal L}g)(x,z) + \lambda(x)\partial_z g(x,z)
\end{equation}
If we consider 
\begin{equation}
w(x,z,t) = {\mathbb E}_{\mu}\!\left[f(X_t)
\exp\!\left(-Z_t\right)\mid X_0=x,Z_0=z
\right]
\end{equation}
Then, with $g(x,z) = f(x)e^{-z}$, we have
\begin{equation}
\partial_t w(x,z,t) = {\mathcal A}w(x,z,t),\quad w(x,z,0) = f(x)e^{-z}
\end{equation}
Now observe that $v(x,t) = w(x,0,t)$, hence 
\begin{equation}
\partial_t v(x,t) = {\mathcal L}v(x,t) +\lambda(x) \partial_z w(x,z,t)\Big|_{z=0}
\end{equation}
Due to 
\begin{align}
\partial_z w(x,z,t)
&=
\mathbb E_{\mu}\!\left[
f(X_t)\,\partial_z \exp\!\left(-\int_0^t\lambda(X_s)\,ds-z\right)
\;\middle|\; X_0=x
\right]\\
&=
-\mathbb E_{\mu}\!\left[
f(X_t)\,\exp\!\left(-\int_0^t\lambda(X_s)\,ds-z\right)
\;\middle|\; X_0=x
\right].
\end{align}
hence 
\begin{equation}
\partial_z w(x,0,t) = -v(x,t).
\end{equation}
and therefore we arrive at the \emph{Feynman-Kac} PDE
\begin{equation}\label{eq:fk}
\partial_t v(x,t) = {\mathcal L}v(x,t) -\lambda(x) v(x,t).
\end{equation}
In particular, with $v(x,0) =1$, this becomes 
\begin{equation}
v(x,t) = {\mathbb E}_{\mu} \left[\exp\left(- \int_0^t\lambda(X_s)ds \right)\middle| X_0 = x \right]
\end{equation}
The problem of computing $\mu_t(f) = {\mathbb E}[v(X_0,t)]$ has thus transformed into a problem of solving a PDE. 




\subsection*{Particle interpretation of Feynman-Kac formula}
In the case of Brownian dynamics we gave an interpretation of the density from a many-particle view. In particular, we viewed the probability density as describing a normalized concentration for a large particle system. The same can be done for the Feynman-Kac PDE, but we now need to view it as relating to densities in a system with changing particle number.

 To see this, consider a population of particles which die at rate $\lambda(X_t)$. Define $N(X_t \in B\mid X_0)$ to be the number of particles alive whose positions are in a measurable set $B$ and started from initial point $X_0$. We then define a number density $n(x,t)$, similarly to a probability density, as 
\begin{equation}
N(X_t \in B\mid X_0) = \int_B n(x,t)\,dx
\end{equation}
Consider how $n(x,t)$ changes due to flux and particle death; by a conservation of mass argument similar to the one we used when deriving the flux for Brownian dynamics, we have 
\begin{equation}
\partial_t n(x,t) = -\nabla \cdot j(x,t) - \lambda(x)n(x,t) = {\mathcal L}^*n(x,t) - \lambda(x)n(x,t).
\end{equation}
We can then define  
\begin{equation}
v_t(f) = \int n(x,t)\,f(x)\,dx 
\end{equation}
Thus $v_t(f)$ has the adjoint equation. We can create a normalized density associated with $n$ to obtain population expectations. This means that another way to obtain $\mu(f)$ is via a population simulation. This is generally more stable than evaluating the expectation directly. 




\end{document}
